\section{Discussion}\label{sec:discussion}

Precision attainment requires both discovering B and reaching the tighter target after discovery. Experiment~1 considers a general run of the compared methods. At \(c=0\), all three adaptive methods have the same discovery count, yet BO-EI and TuRBO-1 convert more discoveries into precise solutions than CMA-ES. As \(c\) increases, both BO variants visit a fixed neighbourhood of \(b\) less frequently. The decline in discovery therefore reflects changed sampling behaviour and inability of the BO algorithm to find the smaller basin.

For BO-EI, a smooth Gaussian process may extrapolate high values from the outer part of B across its unobserved core, leaving higher EI near \(a\).  In Experiment~2, uniform probes imporve coverage directly, whereas increasing \(\xi\) still ranks points using the fitted posterior. The probe policy also changes what the GP learns, so its measured recovery benefit includes subsequent EI proposals. Experiment~3 illustrates how the role of probes changes with the initial information. Once B is observed, the ten probes replace evaluations that EI could devote to refinement, which is reflected in the decrease of precision decrease.

For TuRBO-1, discoveries during restarted designs support the role of renewed global coverage. Faster contraction also changes proposals before a restart and can reduce the time spent refining the current incumbent. 

In CMA-ES, a single candidate can discover basin B and immediately improve the incumbent, while the search distribution updates once per generation. Larger populations draw more candidates per generation and allow fewer updates under a fixed evaluation budget. In Experiment 2, among matched runs that both settled near B, the smaller population generally settled sooner. Across all runs, the larger population finished near B more often. Halving the mean update also improved discovery and precision from near B, consistent with keeping the search centred closer to the global minimiser.  Experiment~3 tests the dispersion explanation more directly. From the same mean near B, a smaller step size concentrates more evaluations in the core and substantially improves precision. From near A, the same concentration favours refinement around the inferior minimiser. The gain from slower mean adaptation near B in Experiment~2 is also consistent with retaining a promising search centre.
